% =============================================================================
\section{Why \bert{} Was a Turning Point}
\label{sec:turning_point}
% =============================================================================

\subsection{Pre-training Objectives}

\bert{} is pre-trained on two self-supervised tasks using only unlabelled text
(Wikipedia + BookCorpus, $\approx$3.3 billion words):

\subsubsection{Masked Language Model (MLM)}

Inspired by the Cloze task~\cite{taylor1953cloze}, 15\% of input tokens are
randomly masked. \bert{} is trained to predict the original token from its
\emph{bidirectional} context:
\begin{equation}
  \mathcal{L}_{\text{MLM}}
    = -\sum_{i \in \mathcal{M}}
      \log P\!\left(x_i \mid x_1, \ldots, x_{i-1},
      \mask{}, x_{i+1}, \ldots, x_n\right)
  \label{eq:mlm}
\end{equation}
where $\mathcal{M}$ is the set of masked positions. Because a correct prediction
requires integrating context from \emph{both} sides, this objective inherently
forces the model to learn deeply bidirectional representations. Traditional
left-to-right language models cannot use this objective without a fundamental
architectural change.

\subsubsection{Next Sentence Prediction (NSP)}

Many NLP tasks (\textit{e.g.}\ question answering and natural language
inference) require understanding the \emph{relationship} between two sentences.
\bert{} is therefore also trained on a binary classification task: given a
sentence pair $(A, B)$, predict whether $B$ is the sentence that actually
follows $A$ in the corpus or a randomly sampled sentence. This trains the model
to encode inter-sentence coherence in the \cls{} token.

\subsection{Fine-tuning Paradigm}

After pre-training, \bert{}'s weights serve as a starting point for any
downstream NLP task. A lightweight task-specific head is appended:

\begin{itemize}[leftmargin=2em, itemsep=2pt]
  \item \textbf{Classification} (sentiment, topic): linear layer on the
        \cls{} representation.
  \item \textbf{Token labelling / NER}: token-level classifier on each output
        position.
  \item \textbf{Extractive QA}: two linear layers predicting span start and end
        positions.
  \item \textbf{Natural Language Inference}: classification on \cls{} for
        entailment / contradiction / neutral.
\end{itemize}

Fine-tuning requires only 2--4 epochs, making \bert{} highly data-efficient.
The same pre-trained model handles all four task types with no structural
changes to the core encoder.

\subsection{Impact: State-of-the-Art Results}

Upon release, \bert{} set new state-of-the-art results across eleven NLP
benchmarks~\cite{devlin2019bert}:

\begin{table}[H]
\centering
\caption{Selected benchmark results at the time of \bert{}'s
  release~\cite{devlin2019bert}.}
\label{tab:results}
\begin{tabular}{llll}
  \toprule
  \textbf{Benchmark} & \textbf{Task} & \textbf{Prior SOTA} & \textbf{\bert{}} \\
  \midrule
  GLUE~\cite{wang2018glue}
    & Multi-task NLU            & 72.8  & 80.4 \\
  SQuAD v1.1~\cite{rajpurkar2016squad}
    & Reading comprehension (F1) & 91.7  & 93.2 \\
  SQuAD v2.0~\cite{rajpurkar2018squad2}
    & Unanswerable QA (F1)       & 78.0  & 83.1 \\
  MultiNLI~\cite{williams2018multinli}
    & Natural language inference & 86.7  & 86.7$^{+}$ \\
  \bottomrule
\end{tabular}
\end{table}

On SQuAD v1.1, \bert{} exceeded human performance on the F1 metric. Crucially,
these results came from a \emph{single} pre-trained model fine-tuned separately
for each task, demonstrating that rich, general-purpose language representations
can be learned once and deployed everywhere.

\subsection{Why It Was a Paradigm Shift}

Four reasons make \bert{} a genuine turning point in NLP:

\begin{enumerate}[leftmargin=2em, itemsep=3pt]
  \item \textbf{Deep bidirectionality at scale.} Unlike ELMo's shallow
        concatenation of separately trained LSTMs, \bert{} integrates left and
        right context jointly in every layer, yielding fundamentally richer
        token representations.

  \item \textbf{Transfer learning for NLP.} \bert{} demonstrated that the
        ``pre-train on unlabelled data, fine-tune on small labelled data''
        paradigm --- long established in computer vision --- works powerfully for
        language. This shifted NLP research away from task-specific architectures
        towards general-purpose pre-trained models.

  \item \textbf{Open release.} Google released \bert{}'s weights, code, and
        pre-training procedure publicly, enabling rapid community adoption and
        spawning a generation of derivative models: RoBERTa~\cite{liu2019roberta},
        DistilBERT~\cite{sanh2019distilbert}, ALBERT~\cite{lan2019albert},
        domain-specific variants, and multilingual extensions.

  \item \textbf{Contextual representations by design.} Every output embedding
        depends on the full sentence context through every attention layer,
        resolving polysemy as an architectural consequence rather than an
        optional post-processing step.
\end{enumerate}
